\documentclass[11pt]{article}
\usepackage[margin=1in]{geometry}
\usepackage{amsmath,amssymb,amsfonts,mathtools}
\usepackage{array,booktabs,enumitem,graphicx,xcolor,hyperref}
\usepackage[T1]{fontenc}
\usepackage[utf8]{inputenc}
\title{The Effective String of 3D Gauge Systems at the Deconfining Transition}
\author{Source-backed arXiv sample}
\date{}
\begin{document}
\maketitle
\section{Extracted Lists}

The list environments below are extracted from arXiv LaTeX source and wrapped for pdfLaTeX validation.

\begin{description}
\item{SU(2)]}
 In this case there is a good convergence among different published
results both for the string tension and for the critical temperature.
Data on $T_c$ are reported in~\cite{tc1} for lattices with length $L=2$
in
the imaginary time direction; in~\cite{tc2} for $L=4$ and in~\cite{tc3}
for $L=5$ and 6. Scaling seems to be fulfilled already for $L>4$
and allows a rather precise estimation: $T_c\beta a=1.4\pm0.1
$~\cite{tc3} (where $a$ is the lattice spacing and $\beta$ the inverse
of the coupling constant).
Data on $\sigma$ are reported in~\cite{sigma1} and more
recently in~\cite{sigma2}, in both case obtained with Wilson loops.
Scaling seems to be fulfilled already for $\beta>6.$ (in agreement with
the $T_c$ behaviour). These two sets of data agree and
give $\sigma(\beta a)^2=2.23\pm0.04$.  This implies
$T_c/\sqrt{\sigma}=0.94\pm0.03$ in agreement with our critical
temperature for the minimal radius $r_1$,
$T_c/{\sqrt\sigma}={\sqrt3}/{\sqrt\pi}=0.98$.
 Indeed this result was
already presented in~\cite{fp} to support the Nambu-Goto string picture, and
has not been significatively changed by the new high precision data of
ref~\cite{sigma2}.

\item{$Z_2$]}
 In this case a rather precise estimate exists for the string
tension ~\cite{noiz2}~:
$\sigma a^2(\beta_c-\beta)^{-2\nu}=3.64\pm0.09$~, with
$\beta_c=0.7614$ and $\nu=0.63$.
For the critical
temperature, the scaling region seems to start from $L=4$. Data obtained
with the Montecarlo renormalization group approach~\cite{wan2} for
$L=4,8$, give $T_c a(\beta_c-\beta)^{-\nu}
=2.3\pm0.1$, while a slightly lower value was
obtained using finite size scaling techniques~\cite{wan1,brower}
for $L=4$:  $T_c a(\beta_c-\beta)^{-\nu}=2.2\pm0.1$. Taking
into account all the uncertainties we
can assume: $T_c/\sqrt{\sigma}=1.17\pm0.1$ which is in acceptable
agreement with our solution.

\end{description}
\end{document}
